\documentclass[12pt,a4paper]{article}

%% Packages
\usepackage{times}
\usepackage[a4paper, left=1.60cm, right=1.60cm, top=1.20cm, bottom=0.60cm]{geometry}
\usepackage{setspace}
\usepackage{graphicx}
\usepackage{booktabs}
\usepackage{array}
\usepackage{hyperref}
\usepackage{enumitem}
\usepackage{titlesec}
\usepackage{caption}
\usepackage{float}
\usepackage{amsmath}

%% Spacing and font
\setstretch{1.15}
\setlength{\parindent}{0pt}
\setlength{\parskip}{6pt}

%% Section heading style
\titleformat{\section}{\bfseries\normalsize}{\thesection.}{0.5em}{}
\titleformat{\subsection}{\bfseries\normalsize}{\thesubsection}{0.5em}{}
\titlespacing*{\section}{0pt}{10pt}{4pt}
\titlespacing*{\subsection}{0pt}{8pt}{3pt}

%% Hyperlink style
\hypersetup{colorlinks=true, linkcolor=black, urlcolor=blue, citecolor=black}

\begin{document}

%% Title
\begin{center}
{\fontsize{24}{28}\bfseries
Failure Intelligence Engine: An Open Source Framework for\\
Real-Time Adversarial Attack Detection and Hallucination\\
Monitoring in Large Language Models}

\vspace{10pt}
{\fontsize{16}{20}\bfseries Ayush Singh}

\vspace{4pt}
{\small B.Tech Student, Department of Computer Science (AI and Data Science),\\
Pimpri Chinchwad University, Pune, India}\\
{\small \href{mailto:ayushsingh355vns@gmail.com}{ayushsingh355vns@gmail.com}}
\end{center}

\vspace{10pt}

%% Abstract
\textbf{Abstract}

Large Language Models (LLMs) are now being used in real-world applications
every day. But they are still easy to attack. Common attacks include prompt
injection, jailbreaking, many-shot conditioning, and model extraction.
At the same time, LLMs often produce wrong or made-up information that cannot
be trusted without checking. Most existing tools either need expensive GPU
hardware, use closed APIs, or only protect against one type of attack.

This paper presents the Failure Intelligence Engine (FIE), an open source
framework that combines a 10-layer adversarial detection pipeline with a
hallucination monitoring system. It can be used as a Python SDK with no GPU
needed. FIE is tested against JailbreakBench [1] and HarmBench [2], two
published evaluation benchmarks. On JailbreakBench, FIE achieves 98.6\%
recall and 97.9\% F1 score, which is better than Meta's Llama Prompt Guard
2-86M (64.9\% recall), while running fully offline. On HarmBench, FIE
achieves 70.6\% recall with 93.4\% precision across 7 harm categories.
All source code and benchmark scripts are publicly available.

\textbf{Keywords:} Large Language Models, Prompt Injection, Jailbreak Detection,
Adversarial Attack, Hallucination Monitoring, AI Safety, LLM Security,
Many-Shot Jailbreaking, Model Extraction

\section{Introduction}

Large Language Models such as GPT-4, Claude, and Llama are now used in
customer support, coding tools, healthcare queries, and education. As more
people use them, protecting them from misuse becomes more important.

Two types of failures matter the most in production. The first is adversarial
attacks, where a user writes a special prompt to trick the model into ignoring
its instructions, revealing its system prompt, or producing harmful content.
The second is hallucination, where the model confidently gives wrong information
even when it does not know the answer [3].

Existing tools do not fully solve this problem. Meta's Llama Prompt Guard [4]
handles some jailbreak patterns but misses semantic attacks. Rebuff only covers
injection patterns. LlamaGuard [11] needs a separate GPU to run. No single open
source tool covers all attack types with production-level security and an easy
deployment path.

FIE was built to fill this gap. It provides:

\begin{itemize}
  \item A 10-layer adversarial detection pipeline covering prompt injection,
        jailbreaks, many-shot conditioning, GCG suffix attacks, indirect injection,
        model extraction, and system prompt leakage.
  \item A hallucination monitoring system using a shadow model jury, an XGBoost
        classifier with AUC-ROC of 0.840, and ground truth verification.
  \item A Python SDK (\texttt{pip install fie-sdk}) that runs fully offline with
        bundled ML models. No GPU, no API key, and no server are required.
  \item A production-ready REST API deployed on Google Cloud Run with automated
        CI/CD, security hardening, and rate limiting.
\end{itemize}

\section{Related Work}

\textbf{Adversarial Attacks on LLMs.}
Prompt injection was formally described by Perez and Ribeiro (2022) [5], showing
that LLMs can be manipulated by embedding instructions inside user input. Zou
et al. (2023) [6] introduced Greedy Coordinate Gradient (GCG) attacks, which
use gradient-based search to find adversarial suffixes that reliably cause
harmful outputs. Chao et al. (2023) [7] introduced PAIR (Prompt Automatic
Iterative Refinement), where an attacker LLM rewrites prompts until a target
model is jailbroken. Anil et al. (2024) [8] from Google DeepMind introduced
many-shot jailbreaking, showing that embedding many scripted Q/A pairs in a
prompt progressively bypasses safety training.

\textbf{Hallucination in LLMs.}
Ji et al. (2023) [3] provide a survey of LLM hallucination, putting it into
categories such as factual inconsistency, fabrication, and reasoning errors.
TruthfulQA [9] and HaluEval [10] are standard benchmarks for measuring how
often models produce false statements.

\textbf{Existing Detection Tools.}
Meta's Llama Prompt Guard 2 [4] is a fine-tuned classifier for jailbreak and
injection detection. It gets 64.9\% recall on JailbreakBench [1] but requires
running a separate language model. Rebuff only covers injection patterns.
LlamaGuard [11] classifies inputs and outputs for policy violations but needs
GPU inference. FIE is different because it combines multiple detection layers
into one offline SDK with no infrastructure cost.

\section{System Architecture}

FIE works in two modes: a local SDK mode and a server pipeline mode.

In local SDK mode, a developer wraps their LLM call with the \texttt{@monitor}
decorator or calls \texttt{scan\_prompt()} directly. All detection runs inside
the application using bundled models. No network call is made.

In server mode, the prompt is also sent to the FIE FastAPI backend, which runs
additional layers including FAISS-based semantic search against 1,000+ labeled
attack vectors, a three-model shadow jury (Llama-3.3-70B, DeepSeek-R1,
Qwen-QWQ-32B via Groq API), and ground truth verification using Wikidata
and Serper search.

\section{Detection Methodology}

\subsection{Adversarial Detection Pipeline}

The pipeline has 10 layers. Each layer runs independently and produces a
confidence score. The final score is the highest confidence across all layers
that fired. Table 1 describes each layer.

\begin{table}[H]
\centering
\caption{Table 1: FIE Adversarial Detection Layers}
\vspace{4pt}
\begin{tabular}{p{1.0cm} p{2.5cm} p{2.5cm} p{6.5cm}}
\toprule
\textbf{Layer} & \textbf{Scope} & \textbf{Method} & \textbf{What It Catches} \\
\midrule
1  & SDK + Server & Regex library    & Direct injection, jailbreak personas, token smuggling, instruction override \\
2  & SDK + Server & PromptGuard      & Keyword combinations with leet-speak normalization \\
3b & SDK + Server & Many-shot detector & Scripted Q/A exchanges conditioning the model \\
4  & SDK + Server & Indirect injection & Attacks embedded in documents or URLs \\
5  & SDK + Server & GCG scanner      & High-entropy adversarial suffixes \\
6  & SDK + Server & Perplexity proxy  & Base64, Caesar cipher, Unicode lookalike attacks \\
7  & SDK + Server & PAIR classifier  & Iteratively rephrased natural language jailbreaks \\
3  & Server only  & FAISS search     & Semantic similarity to 1,000+ labeled attacks \\
8  & Server only  & Consistency check & Output disconnected from the input topic \\
9  & Server only  & LLM intent check  & PAIR-style attacks that pass structural layers \\
\bottomrule
\end{tabular}
\end{table}

\subsection{Many-Shot Jailbreak Detection}

Many-shot jailbreaking [8] works by putting many Human/Assistant exchanges
before a harmful request. The long context makes the model continue responding
cooperatively even when the final question is harmful.

FIE detects this using three signals: (1) the number of Human/Assistant pairs
in the prompt, (2) the fraction of exchanges that mention harmful topics, and
(3) whether the topic moves from harmless to harmful as the conversation goes on.

For 4 to 7 pairs, a harmful signal must be present before an attack is flagged.
This design was necessary to achieve 0\% false positive rate on benign
educational Q/A prompts. For 8 or more pairs, volume alone is treated as
suspicious because the conditioning effect described in [8] becomes strong at
this scale.

\subsection{Model Extraction Detection}

Model extraction attacks involve systematically probing a model to copy its
behavior [12]. FIE tracks three patterns per user session within a one-hour window:
capability probing prompts (e.g., ``what can you do?''), output harvesting using
near-identical prompts detected by Jaccard similarity above 0.85, and high
request rates above 50 requests per hour with systematic prompt variation.

\subsection{Prompt Exfiltration Detection}

Prompt exfiltration happens when a model output reveals its system prompt. FIE
detects this using canary tokens injected into shadow model system prompts, and
six structural regex patterns that match numbered rule lists, role definition
phrases, markdown system-prompt headers, and direct disclosure statements.

\subsection{Hallucination Detection}

Hallucination detection uses four methods on the model output. Shannon entropy
scoring measures how uncertain the model is. Agreement scoring checks whether
three independent shadow models give the same answer. An XGBoost classifier
trained on 2,477 labeled examples predicts hallucination probability. Ground
truth verification cross-checks answers against Wikidata and optionally Serper.

\section{Evaluation}

\subsection{JailbreakBench Evaluation}

JailbreakBench [1] has 282 adversarial prompts across GCG, JBC, and PAIR
attack methods, plus 100 benign prompts from Stanford Alpaca. Table 2 shows
FIE's results compared to Meta's Llama Prompt Guard 2.

\begin{table}[H]
\centering
\caption{Table 2: JailbreakBench Results -- FIE vs Llama Prompt Guard 2}
\vspace{4pt}
\begin{tabular}{lcccccc}
\toprule
\textbf{System} & \textbf{Recall} & \textbf{PAIR} & \textbf{GCG} & \textbf{JBC} & \textbf{FPR} & \textbf{F1} \\
\midrule
FIE v1.4.1 (offline) & \textbf{98.6\%} & \textbf{96.3\%} & \textbf{99.0\%} & \textbf{100\%} & 8.0\% & \textbf{97.9\%} \\
Llama Prompt Guard 2-86M & 64.9\% & 32.9\% & 56.0\% & 100\% & 0.0\% & 78.7\% \\
Llama Prompt Guard 2-22M & 53.5\% & 15.8\% & 38.0\% & 100\% & 1.0\% & 69.6\% \\
\bottomrule
\end{tabular}
\end{table}

FIE achieves 98.6\% overall recall compared to 64.9\% for Meta's best model.
The PAIR recall improvement from FIE v1.1 (3.7\%) to v1.4.1 (96.3\%) is
driven by adding the bundled LinearSVM classifier trained on sentence embeddings.

\subsection{HarmBench Evaluation}

HarmBench [2] covers 320 harmful behaviors across 7 categories plus 200
benign prompts. Table 3 shows FIE's per-category results.

\begin{table}[H]
\centering
\caption{Table 3: HarmBench Per-Category Detection Results}
\vspace{4pt}
\begin{tabular}{lc}
\toprule
\textbf{Harm Category} & \textbf{Detection Rate} \\
\midrule
Harassment and Bullying & 95.2\% \\
Misinformation and Disinformation & 92.6\% \\
Cybercrime and Intrusion & 90.4\% \\
Illegal Activity & 88.7\% \\
Harmful Content & 83.3\% \\
Chemical and Biological & 66.7\% \\
Copyright Violations & 23.8\% \\
\midrule
\textbf{Overall Recall} & \textbf{70.6\%} \\
\textbf{Precision} & \textbf{93.4\%} \\
\textbf{F1} & \textbf{80.4\%} \\
\bottomrule
\end{tabular}
\end{table}

Copyright violations score lowest (23.8\%) because such content does not
contain the injection or jailbreak syntax that FIE's structural layers rely on.

\subsection{New Attack Type Evaluation}

Table 4 shows results for the three detection modules added in v1.4.1.

\begin{table}[H]
\centering
\caption{Table 4: FIE v1.4.1 New Attack Type Evaluation}
\vspace{4pt}
\begin{tabular}{lcccc}
\toprule
\textbf{Attack Type} & \textbf{Recall} & \textbf{FPR} & \textbf{Precision} & \textbf{F1} \\
\midrule
Many-Shot Jailbreak (full pipeline) & 100.0\% & 0.0\% & 100.0\% & 100.0\% \\
Model Extraction & 83.3\% & 0.0\% & 100.0\% & 90.9\% \\
Prompt Exfiltration & 100.0\% & 0.0\% & 100.0\% & 100.0\% \\
\bottomrule
\end{tabular}
\end{table}

\subsection{Hallucination Detection Benchmark}

Table 5 shows hallucination detection results on 2,477 labeled examples from
TruthfulQA [9], HaluEval [10], and MMLU.

\begin{table}[H]
\centering
\caption{Table 5: Hallucination Detection Results (XGBoost v4)}
\vspace{4pt}
\begin{tabular}{lccc}
\toprule
\textbf{Method} & \textbf{Recall} & \textbf{FPR} & \textbf{AUC-ROC} \\
\midrule
POET rule-based baseline & 56.4\% & 38.7\% & -- \\
XGBoost v3 (1,757 examples) & 63.6\% & 38.6\% & 0.677 \\
XGBoost v4 (2,477 examples) & \textbf{68.2\%} & \textbf{8.4\%} & \textbf{0.840} \\
\bottomrule
\end{tabular}
\end{table}

XGBoost v4 reduces the false positive rate from 38.6\% to 8.4\% while
improving recall by 4.6 percentage points over v3, and by 11.8 percentage
points over the rule-based baseline.

\section{Production Deployment}

FIE is deployed on Google Cloud Run in the \texttt{asia-south1} region with
an automated CI/CD pipeline on GitHub Actions. Every push to the main branch
runs secret scanning using Gitleaks, dependency auditing using pip-audit,
static analysis using Ruff, matrix testing across Python 3.10 / 3.11 / 3.12,
Docker image build and push to Google Artifact Registry, and automated
deployment with post-deployment health verification.

Authentication uses Workload Identity Federation, which means no service
account keys are stored anywhere. The API has rate limiting at 100 requests
per minute per IP, security headers (HSTS, Content Security Policy,
X-Frame-Options), and CORS configured with specific allowed origins.

\section{Discussion}

FIE shows that high-recall adversarial detection is possible without GPU
infrastructure. This makes it usable by individual developers and small teams
who cannot afford dedicated model serving.

The main limitation is the 8\% false positive rate on JailbreakBench benign
prompts. This comes mostly from the PAIR classifier flagging borderline prompts.
Future work includes per-layer confidence calibration and testing on multilingual
attack datasets.

The low copyright violation detection (23.8\%) reflects a real gap -- copyright
content does not contain injection patterns. A semantic content classifier
trained specifically on copyright examples would improve this. The model
extraction detector also misses some output-harvesting patterns when Jaccard
similarity falls below the threshold due to minor word-level variation.
Embedding-based similarity would improve recall for this case.

\section{Conclusion}

This paper presented FIE, an open source framework for real-time adversarial
attack detection and hallucination monitoring in large language models. FIE
achieves 98.6\% recall on JailbreakBench, which is 33.7 percentage points
higher than Meta's Llama Prompt Guard 2-86M, while running fully offline
with no GPU.

The system is deployed in production on Google Cloud Run, published as a
Python SDK on PyPI, and tested against two publicly available benchmarks
covering 602 attack prompts and 300 benign prompts across multiple attack
types and harm categories.

FIE is open source under the Apache-2.0 license. The code, SDK, and
benchmark scripts are at:
\url{https://github.com/AyushSingh110/Failure_Intelligence_System}

\section{Acknowledgement}

The author thanks the maintainers of JailbreakBench, HarmBench, TruthfulQA,
and HaluEval for making their datasets publicly available. Benchmark
evaluations used Groq API inference. Claude Code was used as a coding
assistant during development.

\section{References}

\begin{enumerate}[leftmargin=*, label=\arabic*.]

\item Chao P., Robey A., Dobriban E., Hassani H., Pappas G.J., Wong E.,
      ``JailbreakBench: An Open Robustness Benchmark for Jailbreaking Large
      Language Models'', arXiv:2404.01318, 2024.
      \url{https://arxiv.org/abs/2404.01318}

\item Mazeika M., Phan L., Yin X., Zou A., Wang Z., Mu N., Hendrycks D.,
      ``HarmBench: A Standardized Evaluation Framework for Automated Red
      Teaming and Robust Refusal'', arXiv:2402.04249, 2024.
      \url{https://arxiv.org/abs/2402.04249}

\item Ji Z., Lee N., Frieske R., Yu T., Su D., Xu Y., Shi F., Hooi B.,
      ``Survey of Hallucination in Natural Language Generation'',
      ACM Computing Surveys, 2023, 55 (12), 1--38.
      \url{https://arxiv.org/abs/2202.03629}

\item Meta AI, ``Llama Prompt Guard 2'', Meta AI Research, 2024.
      \url{https://ai.meta.com/research/publications/llama-prompt-guard-2/}

\item Perez E., Ribeiro I., ``Ignore Previous Prompt: Attack Techniques
      for Language Models'', arXiv:2211.09527, 2022.
      \url{https://arxiv.org/abs/2211.09527}

\item Zou A., Wang Z., Kolter J.Z., Fredrikson M., ``Universal and
      Transferable Adversarial Attacks on Aligned Language Models'',
      arXiv:2307.15043, 2023.
      \url{https://arxiv.org/abs/2307.15043}

\item Chao P., Robey A., Dobriban E., Hassani H., Pappas G.J., Wong E.,
      ``Jailbreaking Black Box Large Language Models in Twenty Queries'',
      arXiv:2310.08419, 2023.
      \url{https://arxiv.org/abs/2310.08419}

\item Anil C., Durmus E., Sharma M., Benton J., Kundu S., Batson J., et al.,
      ``Many-Shot Jailbreaking'', Anthropic Technical Report, 2024.
      \url{https://www.anthropic.com/research/many-shot-jailbreaking}

\item Lin S., Hilton J., Evans O., ``TruthfulQA: Measuring How Models
      Mimic Human Falsehoods'', Proceedings of ACL, 2022.
      \url{https://arxiv.org/abs/2109.07958}

\item Li J., Cheng X., Zhao W.X., Nie J.Y., Wen J.R., ``HaluEval: A
      Large-Scale Hallucination Evaluation Benchmark for Large Language
      Models'', Proceedings of EMNLP, 2023.
      \url{https://arxiv.org/abs/2305.11747}

\item Inan H., Upasani K., Chi J., Rungta R., Iyer K., Mao Y., et al.,
      ``Llama Guard: LLM-based Input-Output Safeguard for Human-AI
      Conversations'', arXiv:2312.06674, 2023.
      \url{https://arxiv.org/abs/2312.06674}

\item Tramer F., Zhang F., Juels A., Reiter M.K., Ristenpart T.,
      ``Stealing Machine Learning Models via Prediction APIs'',
      Proceedings of USENIX Security Symposium, 2016.
      \url{https://arxiv.org/abs/1609.02943}

\end{enumerate}

\end{document}
